% Lösungen Einheit 3 von 5 – Winkelsummen, Dreiecke und Vierecke (zusammenbau v0.1)
\einheitenkopf{Einheit 3 von 5 · Winkelsummen, Dreiecke und Vierecke}

\erg{16}{a) Dreieck ADC}
\erg{17}{a) rechtwinklig \quad b) $\gamma = 180^\circ - 52^\circ - 71^\circ = 57^\circ$ \quad c) $180^\circ - 51{,}4^\circ - 76{,}9^\circ = 51{,}7^\circ$ \quad d) $180^\circ - 2 \cdot 67^\circ = 46^\circ$ \quad e) $180^\circ : 3 = 60^\circ$ \quad f) $90^\circ - 38^\circ = 52^\circ$ \quad g) $360^\circ - 84^\circ - 97^\circ - 103^\circ = 76^\circ$ \quad h) im Teildreieck ADC: $\gamma_1 = 180^\circ - 90^\circ - 27{,}6^\circ = 62{,}4^\circ$ \quad i) BC = 5 cm (gleiche Winkel bei A und B) \quad j) α und β messen je $45^\circ$, zusammen $90^\circ$; nach der Winkelsumme bleiben für γ $180^\circ - 90^\circ = 90^\circ$. \quad k) gibt es zwei Paare gleich langer Nachbarseiten \quad l) BC = AC = 7 cm (gleiche Winkel bei A und B); $\gamma = 180^\circ - 2 \cdot 66^\circ = 48^\circ$}
\erg{18}{a) 3 Dreiecke: $3 \cdot 180^\circ = 540^\circ$}
\erg{19}{a) Fehler: Die Winkelsumme im Viereck ist $360^\circ$, nicht $180^\circ$; richtig: $360^\circ - 216^\circ = 144^\circ$ \quad b) Die Winkel links und rechts von γ an der Parallelen sind Wechselwinkel zu α und β, also gleich groß. Mit γ bilden sie zusammen einen gestreckten Winkel: $\alpha + \gamma + \beta = 180^\circ$. \quad c) $180^\circ - 2 \cdot 38^\circ = 104^\circ$}
